\begingroup
\setlength{\abovedisplayskip}{3pt plus 1pt minus 1pt}
\setlength{\belowdisplayskip}{3pt plus 1pt minus 1pt}
\setlength{\abovedisplayshortskip}{0pt plus 1pt}
\setlength{\belowdisplayshortskip}{3pt plus 1pt minus 1pt}
\pdfbookmark[2]{Zu IX.}{nexus-remarks-ix}
\begin{center}\large Zu IX.\end{center}

Es sei $p$ eine positive ungerade Zahl, die durch kein Quadrat teilbar ist, und
\[
S_r=\Sigma\left(\frac{s_r}p\right)
\]
wobei $s_r$ alle zu $p$ teilerfremden Zahlen durchlaufen muss, die zwischen $(r-1)\frac p8$ und $r\frac p8$ liegen. Bezeichnet man die Anzahlen der nicht eigentlich äquivalenten, eigentlich primitiven Formen für die Determinanten $-p$ und $-2p$ mit $C_1$ beziehungsweise $C_2$, so ist (vgl. \textsc{Dirichlet}, Recherches sur diverses applications etc., § 11, in \textsc{Crelle}'s Journal XXI)
\[
C_1=2(S_1+S_2),\quad C_2=2(S_1-S_4)
\]
oder
\[
C_1=S_1+S_2+S_3+S_4,\quad C_2=2(S_2+S_3)
\]
je nachdem, ob $p\equiv1$ oder $\equiv3\pmod4$ ist. Ferner bilden die Zahlen $s_1$ und $s_2$ zusammen dieselbe Gesamtheit wie die Zahlen $2s_1$ und $p-2s_4$, ebenso die Zahlen $s_3$ und $s_4$ dieselbe Gesamtheit wie die Zahlen $2s_2$ und $p-2s_3$. Im Fall $p\equiv1\pmod4$ ist die Summe $S_1+S_2+S_3+S_4=0$. Daraus ergeben sich in beiden Fällen zwei weitere Beziehungen zwischen den vier Summen $S_1,S_2,S_3,S_4$, sodass jede durch $C_1$ und $C_2$ ausgedrückt werden kann. Auf diese Weise erhält man, wenn $p\equiv1\pmod4$ ist,
\[
\begin{aligned}
S_1=S_8&=\tfrac14\left(\frac2p\right)C_1+\tfrac14C_2\\
S_2=S_7&=\tfrac14\left(2-\left(\frac2p\right)\right)C_1-\tfrac14C_2\\
S_3=S_6&=-\tfrac14\left(2+\left(\frac2p\right)\right)C_1+\tfrac14C_2\\
S_4=S_5&=\tfrac14\left(\frac2p\right)C_1-\tfrac14C_2
\end{aligned}
\]
\SourceLeaf{312}
und, wenn $p\equiv3\pmod4$ ist,
\[
\begin{aligned}
S_1=-S_8&=\tfrac14\left(3+\left(\frac2p\right)\right)C_1-\tfrac14C_2\\
S_2=-S_7&=-\tfrac14\left(1-\left(\frac2p\right)\right)C_1+\tfrac14C_2\\
S_3=-S_6&=\tfrac14\left(1-\left(\frac2p\right)\right)C_1+\tfrac14C_2\\
S_4=-S_5&=\tfrac14\left(1-\left(\frac2p\right)\right)C_1-\tfrac14C_2
\end{aligned}
\]

Ist $p$ eine Primzahl, so folgen daraus unmittelbar die im Text angegebenen Formeln für die Anzahl der quadratischen Reste in den einzelnen Achteln.
\par\nobreak\begingroup\offinterlineskip\smallskip
\centerline{\rule{35mm}{.35pt}}\smallskip\endgroup
\par\endgroup
